% Content for DWT Discussion subsection
% This file is included in the main conference paper via % Content for DWT Discussion subsection
% This file is included in the main conference paper via % Content for DWT Discussion subsection
% This file is included in the main conference paper via % Content for DWT Discussion subsection
% This file is included in the main conference paper via \input{discussion/DWTDiscussion}

As shown in Figure~\ref{fig:dwt_results} in the Results section, the Discrete Wavelet Transform (DWT) using the Haar wavelet basis exhibits degraded performance on the Texture Images and Icons Images datasets as quantization becomes more aggressive. This degradation in performance can be attributed to the finer details present in these datasets, which are more susceptible to being quantized away during the compression process. The quantization step in the compression pipeline tends to discard high-frequency components that represent fine textures and intricate icon details, leading to a more significant loss of visual information in these particular datasets compared to other image types.

The comparison between Haar and DCT on different image types, as shown in Figures~\ref{fig:dwt_aerial} and~\ref{fig:dwt_kodak}, reveals important characteristics of the Haar wavelet transform. Haar performs better on aerial images because it operates on a global image scale, which is particularly effective for aerial images that typically contain large patches of very similar color, such as oceans or deserts. On the other hand, Haar performs worse on Kodak images since it is less optimized than DCT in the JPEG pipeline, as discussed previously. Additionally, Kodak images generally do not contain large global features of the same color as found in aerial images, which limits the effectiveness of Haar's global approach.


As shown in Figure~\ref{fig:dwt_results} in the Results section, the Discrete Wavelet Transform (DWT) using the Haar wavelet basis exhibits degraded performance on the Texture Images and Icons Images datasets as quantization becomes more aggressive. This degradation in performance can be attributed to the finer details present in these datasets, which are more susceptible to being quantized away during the compression process. The quantization step in the compression pipeline tends to discard high-frequency components that represent fine textures and intricate icon details, leading to a more significant loss of visual information in these particular datasets compared to other image types.

The comparison between Haar and DCT on different image types, as shown in Figures~\ref{fig:dwt_aerial} and~\ref{fig:dwt_kodak}, reveals important characteristics of the Haar wavelet transform. Haar performs better on aerial images because it operates on a global image scale, which is particularly effective for aerial images that typically contain large patches of very similar color, such as oceans or deserts. On the other hand, Haar performs worse on Kodak images since it is less optimized than DCT in the JPEG pipeline, as discussed previously. Additionally, Kodak images generally do not contain large global features of the same color as found in aerial images, which limits the effectiveness of Haar's global approach.


As shown in Figure~\ref{fig:dwt_results} in the Results section, the Discrete Wavelet Transform (DWT) using the Haar wavelet basis exhibits degraded performance on the Texture Images and Icons Images datasets as quantization becomes more aggressive. This degradation in performance can be attributed to the finer details present in these datasets, which are more susceptible to being quantized away during the compression process. The quantization step in the compression pipeline tends to discard high-frequency components that represent fine textures and intricate icon details, leading to a more significant loss of visual information in these particular datasets compared to other image types.

The comparison between Haar and DCT on different image types, as shown in Figures~\ref{fig:dwt_aerial} and~\ref{fig:dwt_kodak}, reveals important characteristics of the Haar wavelet transform. Haar performs better on aerial images because it operates on a global image scale, which is particularly effective for aerial images that typically contain large patches of very similar color, such as oceans or deserts. On the other hand, Haar performs worse on Kodak images since it is less optimized than DCT in the JPEG pipeline, as discussed previously. Additionally, Kodak images generally do not contain large global features of the same color as found in aerial images, which limits the effectiveness of Haar's global approach.


As shown in Figure~\ref{fig:dwt_results} in the Results section, the Discrete Wavelet Transform (DWT) using the Haar wavelet basis exhibits degraded performance on the Texture Images and Icons Images datasets as quantization becomes more aggressive. This degradation in performance can be attributed to the finer details present in these datasets, which are more susceptible to being quantized away during the compression process. The quantization step in the compression pipeline tends to discard high-frequency components that represent fine textures and intricate icon details, leading to a more significant loss of visual information in these particular datasets compared to other image types.

The comparison between Haar and DCT on different image types, as shown in Figures~\ref{fig:dwt_aerial} and~\ref{fig:dwt_kodak}, reveals important characteristics of the Haar wavelet transform. Haar performs better on aerial images because it operates on a global image scale, which is particularly effective for aerial images that typically contain large patches of very similar color, such as oceans or deserts. On the other hand, Haar performs worse on Kodak images since it is less optimized than DCT in the JPEG pipeline, as discussed previously. Additionally, Kodak images generally do not contain large global features of the same color as found in aerial images, which limits the effectiveness of Haar's global approach.
